\documentclass[10pt]{article}

% ---------- Document Identity (update these only) ----------
\newcommand{\DocStage}{}
\newcommand{\DocTitle}{Your Road to Tier One SOF}
\newcommand{\ProgramName}{182-Week Civilian-to-Selection Readiness Pipeline}
\newcommand{\ProgramSubtitle}{}

% ---------- Page ----------
\usepackage[legalpaper,left=0.5in,right=0.5in,top=0.6in,bottom=0.45in]{geometry}

% ---------- Typography ----------
\usepackage[T1]{fontenc}
\usepackage[utf8]{inputenc}
\usepackage{libertinus}
\usepackage{microtype}
\usepackage{parskip}
\usepackage{ragged2e}
\usepackage{textcomp}
\AtBeginDocument{\color{black}}
\AtBeginDocument{\pagecolor{white}}
\AtBeginDocument{\justifying}

% ---------- Landscape pages ----------
\usepackage{pdflscape}

% ---------- Color ----------
\usepackage[table]{xcolor}
\definecolor{StageBlue}{HTML}{1F4E79}
\definecolor{StageBlueDark}{HTML}{163A5A}
\definecolor{LightGray}{HTML}{F2F4F7}
\definecolor{MidGray}{HTML}{D9DEE5}
\definecolor{RuleGray}{HTML}{C7CED8}
\definecolor{HeaderInk}{HTML}{000000}
\definecolor{Ink}{HTML}{000000}

% ---------- Utility ----------
\usepackage{enumitem}
\sloppy
\setlength{\emergencystretch}{2em}

% ---------- Boxes / UI ----------
\usepackage[most]{tcolorbox}

\tcbset{
  charterTitle/.style={
    enhanced,
    colback=StageBlue!4,
    colframe=StageBlue,
    boxrule=1.25pt,
    arc=3pt,
    left=16pt,right=16pt,top=14pt,bottom=14pt,
    borderline north={3.2pt}{0pt}{StageBlue},
    borderline west={4pt}{0pt}{StageBlue},
    borderline south={1.25pt}{0pt}{StageBlue},
    drop shadow={black!25!white},
  },
  charterRule/.style={
    enhanced,
    colback=LightGray,
    colframe=RuleGray,
    boxrule=0.85pt,
    arc=3pt,
    left=12pt,right=12pt,top=10pt,bottom=10pt,
    borderline west={3pt}{0pt}{StageBlue},
  },
  charterBadge/.style={
    enhanced,
    colback=StageBlue,
    coltext=white,
    colframe=StageBlueDark,
    boxrule=0.6pt,
    arc=2pt,
    left=6pt,right=6pt,top=3pt,bottom=3pt,
  }
}
\newtcbox{\Badge}{nobeforeafter,charterBadge}

% ---------- Lists ----------
\setlist[itemize]{leftmargin=1.5em, itemsep=0.25em, topsep=0.25em}
\setlist[enumerate]{leftmargin=1.8em, itemsep=0.25em, topsep=0.25em}

% ---------- TOC ----------
\usepackage{tocloft}
\setcounter{tocdepth}{3}
\setlength{\cftbeforesecskip}{3pt}
\setlength{\cftbeforesubsecskip}{2pt}
\setlength{\cftbeforesubsubsecskip}{0pt}
\renewcommand{\cftsecfont}{\bfseries}
\renewcommand{\cftsecpagefont}{\bfseries}
\renewcommand{\cftsubsecfont}{\normalfont}
\renewcommand{\cftsubsecpagefont}{\normalfont}
\renewcommand{\cftsubsubsecfont}{\normalfont}
\renewcommand{\cftsubsubsecpagefont}{\normalfont}
\setlength{\cftsecindent}{0pt}
\setlength{\cftsubsecindent}{1.25em}
\setlength{\cftsubsubsecindent}{2.8em}
\setlength{\cftsecnumwidth}{2.1em}
\setlength{\cftsubsecnumwidth}{2.8em}
\setlength{\cftsubsubsecnumwidth}{3.6em}

% Remove the built-in TOC heading so only the custom heading is shown
\makeatletter
\renewcommand{\tableofcontents}{\@starttoc{toc}}
\makeatother

% ---------- Header/Footer: page number only ----------
\usepackage{fancyhdr}
\pagestyle{fancy}
\fancyhf{}
\renewcommand{\headrulewidth}{0pt}
\renewcommand{\footrulewidth}{0pt}
\fancyfoot[C]{\thepage}

% ---------- Section formatting ----------
\usepackage{titlesec}

\titleformat{\section}
  {\Large\bfseries\color{Ink}}
  {\thesection}{0.6em}{}
  [\vspace{0.25em}\textcolor{StageBlue}{\rule{\linewidth}{0.8pt}}]

\titleformat{\subsection}
  {\large\bfseries\color{Ink}}
  {\thesubsection}{0.6em}{}

\titleformat{\subsubsection}
  {\normalsize\bfseries\color{Ink}}
  {\thesubsubsection}{0.6em}{}

\titlespacing*{\section}{0pt}{0.95em}{0.35em}
\titlespacing*{\subsection}{0pt}{0.65em}{0.25em}
\titlespacing*{\subsubsection}{0pt}{0.55em}{0.2em}

% ---------- Table engine ----------
\usepackage{tabularray}
\UseTblrLibrary{booktabs}

% ---------- tabularray: GLOBAL table treatment ----------
\SetTblrInner{
  cells = {fg=black, bg=white, valign=t},
}

% ---------- Header helpers ----------
\newcommand{\Hdr}[1]{\shortstack[c]{\strut #1\strut}}

% ---------- Lines ----------
\newcommand{\TitleLine}{\textcolor{StageBlue}{\rule{0.98\linewidth}{0.95pt}}}
\newcommand{\SubTitleLine}{\textcolor{RuleGray}{\rule{0.98\linewidth}{0.65pt}}}

% ---------- Continued on next page ----------
\newcommand{\ContinuedOnNextPageInline}{%
  \par
  \vfill
  \noindent\textcolor{RuleGray}{\rule{\linewidth}{1pt}}\vspace{-2pt}
  \noindent\hfill\textit{Continued on next page}\par\vspace{-10pt}
  \noindent\textcolor{RuleGray}{\rule{\linewidth}{1pt}}
  \clearpage
}

% ---------- PDF hyperlinks ----------
\usepackage[hidelinks]{hyperref}
\hypersetup{
  colorlinks=false,
  pdfborder={0 0 0},
  linktoc=all,
  pdftitle={\DocTitle},
  pdfauthor={\ProgramName}
}

% =========================================================
\begin{document}

\begin{landscape}

\section*{Stage 6.C — Phase-Specific Acquisition Lists Delta-Based}
\phantomsection
\addcontentsline{toc}{section}{Stage 6.C — Phase-Specific Acquisition Lists Delta-Based}

\subsection*{6.C.1 Priority Scale}
\phantomsection
\addcontentsline{toc}{subsection}{6.C.1 Priority Scale}

\begin{itemize}
\item High = required by phase entry or mid-phase window to execute safely and validly; cannot proceed without it.
\item Med = required by end-phase gate window to execute safely and validly; can be phased within the phase but must be present by end gate.
\item Low = optional or convenience upgrade; not required for validity or safety; defer if budget/space constrained.
\end{itemize}

\subsection*{6.C.2 Phase Delta Tables}
\phantomsection
\addcontentsline{toc}{subsection}{6.C.2 Phase Delta Tables}

\subsubsection*{Phase 00 — Bodily Reactivation and Baseline Creation}
\phantomsection

\begingroup
\scriptsize
\setlength{\tabcolsep}{2pt}
\renewcommand{\arraystretch}{1.12}

\begin{longtblr}{
  width=\linewidth,
  rowhead=1,
  colspec = {
    X[l,1.95]
    X[l,0.90]
    X[l,2.40]
    X[l,0.85]
    X[l,2.15]
  },
  hlines,
  vlines,
  cells = {hsep=6pt, vsep=6pt, halign=l, valign=t},
  cell{1-Z}{1-5} = {fg=black},
  row{odd} = {bg=white},
  row{even} = {bg=LightGray},
  row{1} = {bg=StageBlue!15, fg=black, font=\bfseries, halign=l, valign=t},
}
\Hdr{Item or Category Change} &
\Hdr{From Tier → To Tier} &
\Hdr{Why Introduced Here \\ (Training/Test Need)} &
\Hdr{Priority \\ (High/Med/Low)} &
\Hdr{Space / Safety / Budget Notes} \\

7.11 Footwear Systems and Foot Care (supportive shoes; sock system; blister prevention supplies) &
T0 → T1 &
Phase 00 locomotion tolerance and durability stability require a controlled footwear interface; prevents hotspot-to-blister drift that would invalidate execution and gate readiness. &
High &
Space: Minimal. Safety: Controls gait drift and skin breakdown; required for safe step-based Cardio continuity. Budget: Prioritize as first acquisition; do not “make do” with unstable footwear. \\

7.18 Testing, Timing, Measurement, and Course-Setup Tools (timer method; scale; tape; basic measurement posture) &
T0 → T1 &
Phase 00 requires auditable logs and repeatable measurement capture; validity depends on stable timing and basic measurement capability. &
High &
Space: Minimal. Safety: Prevents invalid evidence due to missing tools; supports conservative decisions. Budget: Tier 1 only; no upgrades. \\

7.12 Recovery, Mobility, and Tissue-Care Implements (safe surface/mat; basic mobility aid) &
T0 → T1 &
Phase 00 movement quality and durability require daily low-risk mobility/recovery capability without escalation. &
High &
Space: Minimal. Safety: Reduces predictable stiffness drift and compensations; supports stop-posture compliance. Budget: Tier 1 only; generic item types only. \\

7.13 Nutrition Preparation and Hydration Systems (hydration container; basic carry) &
T0 → T1 &
Hydration feasibility supports safe Cardio continuity and reduces heat-related safety risk; supports adherence and recovery stability. &
High &
Space: Minimal. Safety: Supports conservative Cardio execution; reduces avoidable fatigue drift. Budget: Tier 1 only; no elaboration. \\

7.10 Logistics, Maintenance, Storage, and Sustainment Equipment (drying/storage posture for footwear and basics) &
T0 → T1 &
Prevents repeated execution failure from wet footwear, poor storage, and friction that drives missed sessions and foot breakdown. &
Med &
Space: Low-to-moderate. Safety: Protects feet/skin integrity and reduces illness risk. Budget: Tier 1 only; prioritize drying and storage essentials before convenience upgrades. \\

7.20 Personal Hygiene, Health, and Sanitation Items (hygiene/drying/laundry posture) &
T0 → T1 &
Supports skin integrity and illness-risk reduction under fixed 6-sessions/week cadence; prevents avoidable breakdown that would force \textbf{HOLD} posture. &
Med &
Space: Minimal. Safety: Prevents skin/foot degradation and illness-driven drift. Budget: Essentials only; no expansion. \\

7.16 Operational Self-Management Systems (single system-of-record; planning/journaling) &
T0 → T1 &
Phase 00 requires stable logging and adherence discipline; without it, evidence is inadmissible and gates default \textbf{HOLD}. &
High &
Space: Minimal. Safety: Prevents silent drift and backfilled logs. Budget: Minimal; avoid complexity creep. \\

7.14 Environmental Apparel — Cold Hot Wet Inclement Weather (basic weather-appropriate layers) &
T0 → T1 &
Grand Forks climate volatility threatens continuity and safety; basic apparel prevents unsafe exposure and repeated reslots. &
Med &
Space: Moderate. Safety: Enables safe outdoor exposure when appropriate; otherwise supports indoor substitution posture. Budget: Essentials only; defer expansion. \\

7.15 General Training and Daily Wear Apparel (basic training clothing posture) &
T0 → T1 &
Prevents friction-related noncompliance; supports safe movement and comfort under frequent sessions. &
Low &
Space: Moderate. Safety: Reduces chafing/skin irritation risk. Budget: Minimal; use existing household options when feasible. \\

\end{longtblr}

\endgroup

\vspace{0.65em}
\newpage

\subsubsection*{Phase 01 — Base Conditioning and Hypertrophy}
\phantomsection

\begingroup
\scriptsize
\setlength{\tabcolsep}{2pt}
\renewcommand{\arraystretch}{1.12}

\begin{longtblr}{
  width=\linewidth,
  rowhead=1,
  colspec = {
    X[l,1.95]
    X[l,0.90]
    X[l,2.40]
    X[l,0.85]
    X[l,2.15]
  },
  hlines,
  vlines,
  cells = {hsep=6pt, vsep=6pt, halign=l, valign=t},
  cell{1-Z}{1-5} = {fg=black},
  row{odd} = {bg=white},
  row{even} = {bg=LightGray},
  row{1} = {bg=StageBlue!15, fg=black, font=\bfseries, halign=l, valign=t},
}
\Hdr{Item or Category Change} &
\Hdr{From Tier → To Tier} &
\Hdr{Why Introduced Here \\ (Training/Test Need)} &
\Hdr{Priority \\ (High/Med/Low)} &
\Hdr{Space / Safety / Budget Notes} \\

7.01 Strength and Resistance Training Equipment (minimum safe resistance capability) &
T0 → T1 &
Phase 01 introduces structured resistance exposure as a minimum deployable capability; without it, phase intent cannot be executed safely or consistently. &
Med &
Space: Low-to-moderate. Safety: Enables controlled \textbf{STR} exposure; prevents improvised unsafe substitutes. Budget: Tier 1 only; defer upgrades. \\

7.08 Conditioning, Endurance, and Aerobic Training Tools (indoor substitute capability) &
T0 → T1 &
Phase 01 requires year-round Cardio continuity and safe indoor substitution when outdoor conditions invalidate execution; preserves validity and adherence. &
Med &
Space: Moderate. Safety: Reduces ice/heat/air-quality exposure risk. Budget: Minimal tool posture; defer expansion. \\

\end{longtblr}

\endgroup

Phase 01 is locked to the 24-week variant; duration and window doctrine are reconciled here and are not left as an open Stage 8 issue.

\vspace{0.65em}

\subsubsection*{Phase 02 — Maximal Strength Development}
\phantomsection

\begingroup
\scriptsize
\setlength{\tabcolsep}{2pt}
\renewcommand{\arraystretch}{1.12}

\begin{longtblr}{
  width=\linewidth,
  rowhead=1,
  colspec = {
    X[l,1.95]
    X[l,0.90]
    X[l,2.40]
    X[l,0.85]
    X[l,2.15]
  },
  hlines,
  vlines,
  cells = {hsep=6pt, vsep=6pt, halign=l, valign=t},
  cell{1-Z}{1-5} = {fg=black},
  row{odd} = {bg=white},
  row{even} = {bg=LightGray},
  row{1} = {bg=StageBlue!15, fg=black, font=\bfseries, halign=l, valign=t},
}
\Hdr{Item or Category Change} &
\Hdr{From Tier → To Tier} &
\Hdr{Why Introduced Here \\ (Training/Test Need)} &
\Hdr{Priority \\ (High/Med/Low)} &
\Hdr{Space / Safety / Budget Notes} \\

7.01 Strength and Resistance Training Equipment (barbell/rack-level capability allowed as minimum deployable) &
T1 → T1 (Operational delta: capability class shift) &
Phase 02 may require barbell/rack-level capability as the minimum safe means to execute \textbf{STR} proxy intent; absence forces unsafe improvisation or invalid comparability. &
High &
Space: Moderate-to-high. Safety: Enables controlled load handling; prevents unstable substitutes. Budget: If required by phase card, it becomes functional Tier 1; stage acquisitions ahead of protected windows. \\

\end{longtblr}

\endgroup

\vspace{0.65em}

\subsubsection*{Phase 03 — Converting Strength to Power and Endurance for SOF Selection}
\phantomsection

\begingroup
\scriptsize
\setlength{\tabcolsep}{2pt}
\renewcommand{\arraystretch}{1.12}

\begin{longtblr}{
  width=\linewidth,
  rowhead=1,
  colspec = {
    X[l,1.95]
    X[l,0.90]
    X[l,2.40]
    X[l,0.85]
    X[l,2.15]
  },
  hlines,
  vlines,
  cells = {hsep=6pt, vsep=6pt, halign=l, valign=t},
  cell{1-Z}{1-5} = {fg=black},
  row{odd} = {bg=white},
  row{even} = {bg=LightGray},
  row{1} = {bg=StageBlue!15, fg=black, font=\bfseries, halign=l, valign=t},
}
\Hdr{Item or Category Change} &
\Hdr{From Tier → To Tier} &
\Hdr{Why Introduced Here \\ (Training/Test Need)} &
\Hdr{Priority \\ (High/Med/Low)} &
\Hdr{Space / Safety / Budget Notes} \\

7.07 Load-Bearing, Rucking, and Carry Systems (first minimum deployable ruck capability) &
T0 → T1 (End) &
Locked first-introduction point: Phase 03 is the first phase where minimum deployable ruck/carry capability must exist by end gate to prepare for later load tolerance demands. &
Med &
Space: Moderate. Safety: Requires stable fit/configuration and load verification posture; avoid first-time interface changes inside protected windows. Budget: Essentials only; defer upgrades. \\

7.18 Testing, Timing, Measurement, and Course-Setup Tools (higher complexity gate-day readiness) &
T1 → T2 (End) &
Phase complexity and gate days increase; Tier 2 measurement/posture reduces validity breaks from tool drift and improves reproducibility for multi-domain days. &
Med &
Space: Low. Safety: Reduces invalid evidence risk due to missing/unstable measurement methods. Budget: Tier 2 becomes functional minimum if phase card requires it for admissible testing. \\

7.19 Protective and Assistive Equipment (first minimum deployable protective posture) &
T0 → T1 (End) &
First meaningful point where protective/assistive capability reduces predictable breakdown as load/interface complexity increases. &
Med &
Space: Low. Safety: Supports conservative risk mitigation; prevents avoidable overuse drift. Budget: Essentials only. \\

\end{longtblr}

\endgroup

\vspace{0.65em}

\subsubsection*{Phase 04 — Strength-Endurance and Work Capacity Development}
\phantomsection

\begingroup
\scriptsize
\setlength{\tabcolsep}{2pt}
\renewcommand{\arraystretch}{1.12}

\begin{longtblr}{
  width=\linewidth,
  rowhead=1,
  colspec = {
    X[l,1.95]
    X[l,0.90]
    X[l,2.40]
    X[l,0.85]
    X[l,2.15]
  },
  hlines,
  vlines,
  cells = {hsep=6pt, vsep=6pt, halign=l, valign=t},
  cell{1-Z}{1-5} = {fg=black},
  row{odd} = {bg=white},
  row{even} = {bg=LightGray},
  row{1} = {bg=StageBlue!15, fg=black, font=\bfseries, halign=l, valign=t},
}
\Hdr{Item or Category Change} &
\Hdr{From Tier → To Tier} &
\Hdr{Why Introduced Here \\ (Training/Test Need)} &
\Hdr{Priority \\ (High/Med/Low)} &
\Hdr{Space / Safety / Budget Notes} \\

7.17 Aquatic Training and Water Safety Equipment (first minimum deployable aquatic capability) &
T0 → T1 (End) &
Locked first-introduction point: surface swim becomes phase-relevant and requires verified safety/logistics posture; underwater is not implied. &
Med &
Space: Low-to-moderate. Safety: Supports lawful, safe pool execution and valid logging; no underwater/breath-hold governance implied. Budget: Essentials only; defer upgrades. \\

\end{longtblr}

\endgroup

\vspace{0.65em}

\subsubsection*{Phase 05 — Aerobic Base and Extended Stamina Development}
\phantomsection

\begingroup
\scriptsize
\setlength{\tabcolsep}{2pt}
\renewcommand{\arraystretch}{1.12}

\begin{longtblr}{
  width=\linewidth,
  rowhead=1,
  colspec = {
    X[l,1.95]
    X[l,0.90]
    X[l,2.40]
    X[l,0.85]
    X[l,2.15]
  },
  hlines,
  vlines,
  cells = {hsep=6pt, vsep=6pt, halign=l, valign=t},
  cell{1-Z}{1-5} = {fg=black},
  row{odd} = {bg=white},
  row{even} = {bg=LightGray},
  row{1} = {bg=StageBlue!15, fg=black, font=\bfseries, halign=l, valign=t},
}
\Hdr{Item or Category Change} &
\Hdr{From Tier → To Tier} &
\Hdr{Why Introduced Here \\ (Training/Test Need)} &
\Hdr{Priority \\ (High/Med/Low)} &
\Hdr{Space / Safety / Budget Notes} \\

7.18 Testing, Timing, Measurement, and Course-Setup Tools (extended endurance comparability support) &
T2 → T2 (Operational delta: extend validity controls) &
Extended duration gates amplify route/tool drift; operational emphasis shifts to tighter comparability discipline for longer events. &
Med &
Space: Low. Safety: Reduces invalid test risk; supports reslot decisions when conditions drift. Budget: No upgrades required beyond Tier 2 if already met. \\

\end{longtblr}

\endgroup

\vspace{0.65em}

\subsubsection*{Phase 06 — Lactate Threshold and Power-Endurance Development}
\phantomsection

\begingroup
\scriptsize
\setlength{\tabcolsep}{2pt}
\renewcommand{\arraystretch}{1.12}

\begin{longtblr}{
  width=\linewidth,
  rowhead=1,
  colspec = {
    X[l,1.95]
    X[l,0.90]
    X[l,2.40]
    X[l,0.85]
    X[l,2.15]
  },
  hlines,
  vlines,
  cells = {hsep=6pt, vsep=6pt, halign=l, valign=t},
  cell{1-Z}{1-5} = {fg=black},
  row{odd} = {bg=white},
  row{even} = {bg=LightGray},
  row{1} = {bg=StageBlue!15, fg=black, font=\bfseries, halign=l, valign=t},
}
\Hdr{Item or Category Change} &
\Hdr{From Tier → To Tier} &
\Hdr{Why Introduced Here \\ (Training/Test Need)} &
\Hdr{Priority \\ (High/Med/Low)} &
\Hdr{Space / Safety / Budget Notes} \\

7.14 Environmental Apparel — Cold, Hot, Wet, Inclement Weather (first climate-driven expansion) &
T1 → T2 (End) &
First major point where climate volatility becomes a recurrent validity/safety constraint rather than an occasional inconvenience; expanded apparel reduces forced reslots and unsafe exposures. &
Med &
Space: Moderate. Safety: Enables safer execution in variable conditions; still defaults to indoor substitution when unsafe. Budget: Expand only as needed to prevent repeated validity failures. \\

7.11 Footwear Systems and Foot Care (expanded foot interface redundancy) &
T1 → T2 (End) &
Increased consequence of \textbf{RUN}/\textbf{RUCK} domains requires more robust footcare and interface stability to prevent predictable breakdown. &
Med &
Space: Low. Safety: Reduces hotspot/blister cascade; supports gait stability. Budget: Expand only to eliminate recurrent failure modes. \\

7.12 Recovery, Mobility, and Tissue-Care Implements (expanded recovery tooling) &
T1 → T2 (End) &
Higher consequence and density require stronger recovery capability to prevent injury drift and invalid test outcomes. &
Med &
Space: Low. Safety: Supports durability controls and conservative downshift posture. Budget: Expand only to prevent repeated breakdown. \\

7.10 Logistics, Maintenance, Storage, and Sustainment Equipment (expanded sustainment) &
T1 → T2 (End) &
Increased workload magnifies friction costs; sustainment reduces recurring missed sessions and equipment failures. &
Med &
Space: Moderate. Safety: Prevents wet-footwear and skin-breakdown loops. Budget: Expand only to stabilize continuity. \\

7.07 Load-Bearing, Rucking, and Carry Systems (expanded stability) &
T1 → T2 (End) &
Load tolerance becomes more consequential; expanded stability posture reduces risk of fit drift and load movement that corrupts evidence and safety. &
Med &
Space: Moderate. Safety: Improves load stability; prevents chafing and hotspots. Budget: Expand only for stability, not performance. \\

7.18 Testing, Timing, Measurement, and Course-Setup Tools (maintain Tier 2) &
T2 → T2 (Operational delta: gate-day complexity increases) &
Testing complexity rises; maintain Tier 2 measurement posture to prevent tool drift. &
High &
Space: Low. Safety: Gate evidence invalidity triggers \textbf{HOLD}/reslot; keep tools stable. Budget: Protect this category from drift. \\

\end{longtblr}

\endgroup

\vspace{0.65em}

\subsubsection*{Phase 07 — Advanced Hybrid Overload}
\phantomsection

\begingroup
\scriptsize
\setlength{\tabcolsep}{2pt}
\renewcommand{\arraystretch}{1.12}

\begin{longtblr}{
  width=\linewidth,
  rowhead=1,
  colspec = {
    X[l,1.95]
    X[l,0.90]
    X[l,2.40]
    X[l,0.85]
    X[l,2.15]
  },
  hlines,
  vlines,
  cells = {hsep=6pt, vsep=6pt, halign=l, valign=t},
  cell{1-Z}{1-5} = {fg=black},
  row{odd} = {bg=white},
  row{even} = {bg=LightGray},
  row{1} = {bg=StageBlue!15, fg=black, font=\bfseries, halign=l, valign=t},
}
\Hdr{Item or Category Change} &
\Hdr{From Tier → To Tier} &
\Hdr{Why Introduced Here \\ (Training/Test Need)} &
\Hdr{Priority \\ (High/Med/Low)} &
\Hdr{Space / Safety / Budget Notes} \\

7.16 Operational Self-Management Systems (expanded execution control) &
T1 → T2 (End) &
Hybrid density increases scheduling drift risk; expanded self-management reduces missed sessions, undocumented substitutions, and gate-week chaos. &
Med &
Space: Low. Safety: Prevents unsafe stacking and validity contamination. Budget: Expand only to eliminate recurring drift. \\

\end{longtblr}

\endgroup

\vspace{0.65em}

\subsubsection*{Phase 08 — Structural Durability and Load Tolerance}
\phantomsection

\begingroup
\scriptsize
\setlength{\tabcolsep}{2pt}
\renewcommand{\arraystretch}{1.12}

\begin{longtblr}{
  width=\linewidth,
  rowhead=1,
  colspec = {
    X[l,1.95]
    X[l,0.90]
    X[l,2.40]
    X[l,0.85]
    X[l,2.15]
  },
  hlines,
  vlines,
  cells = {hsep=6pt, vsep=6pt, halign=l, valign=t},
  cell{1-Z}{1-5} = {fg=black},
  row{odd} = {bg=white},
  row{even} = {bg=LightGray},
  row{1} = {bg=StageBlue!15, fg=black, font=\bfseries, halign=l, valign=t},
}
\Hdr{Item or Category Change} &
\Hdr{From Tier → To Tier} &
\Hdr{Why Introduced Here \\ (Training/Test Need)} &
\Hdr{Priority \\ (High/Med/Low)} &
\Hdr{Space / Safety / Budget Notes} \\

7.07 Load-Bearing, Rucking, and Carry Systems (expanded load tolerance readiness) &
T2 → T2 (Operational delta: long-load stability becomes dominant) &
Load tolerance milestones demand stable interface and configuration repeatability to protect feet and validity. &
High &
Space: Moderate. Safety: Long-load failure risk is high; avoid late changes near protected windows. Budget: No upgrades unless required to prevent repeated failure. \\

7.19 Protective and Assistive Equipment (expanded risk control) &
T1 → T2 (End) &
Higher consequence phases require stronger protective posture to prevent predictable breakdown. &
Med &
Space: Low. Safety: Reduces overuse drift and protects joints/skin. Budget: Expand only when needed for stability. \\

\end{longtblr}

\endgroup

\vspace{0.65em}

\subsubsection*{Phase 09 — High-Volume Calisthenics and Stamina}
\phantomsection

\begingroup
\scriptsize
\setlength{\tabcolsep}{2pt}
\renewcommand{\arraystretch}{1.12}

\begin{longtblr}{
  width=\linewidth,
  rowhead=1,
  colspec = {
    X[l,1.95]
    X[l,0.90]
    X[l,2.40]
    X[l,0.85]
    X[l,2.15]
  },
  hlines,
  vlines,
  cells = {hsep=6pt, vsep=6pt, halign=l, valign=t},
  cell{1-Z}{1-5} = {fg=black},
  row{odd} = {bg=white},
  row{even} = {bg=LightGray},
  row{1} = {bg=StageBlue!15, fg=black, font=\bfseries, halign=l, valign=t},
}
\Hdr{Item or Category Change} &
\Hdr{From Tier → To Tier} &
\Hdr{Why Introduced Here \\ (Training/Test Need)} &
\Hdr{Priority \\ (High/Med/Low)} &
\Hdr{Space / Safety / Budget Notes} \\

7.12 Recovery, Mobility, and Tissue-Care Implements (maintain Tier 2) &
T2 → T2 (Operational delta: volume-induced friction control) &
High \textbf{CAL} volume requires consistent recovery tooling to avoid overuse drift and invalid evidence. &
Med &
Space: Low. Safety: Prevents tendon/hand/skin breakdown; supports repeatability. Budget: Maintain; do not overbuild unless failure patterns emerge. \\

\end{longtblr}

\endgroup

\vspace{0.65em}

\subsubsection*{Phase 10 — Advanced Tactical Readiness}
\phantomsection

\begingroup
\scriptsize
\setlength{\tabcolsep}{2pt}
\renewcommand{\arraystretch}{1.12}

\begin{longtblr}{
  width=\linewidth,
  rowhead=1,
  colspec = {
    X[l,1.95]
    X[l,0.90]
    X[l,2.40]
    X[l,0.85]
    X[l,2.15]
  },
  hlines,
  vlines,
  cells = {hsep=6pt, vsep=6pt, halign=l, valign=t},
  cell{1-Z}{1-5} = {fg=black},
  row{odd} = {bg=white},
  row{even} = {bg=LightGray},
  row{1} = {bg=StageBlue!15, fg=black, font=\bfseries, halign=l, valign=t},
}
\Hdr{Item or Category Change} &
\Hdr{From Tier → To Tier} &
\Hdr{Why Introduced Here \\ (Training/Test Need)} &
\Hdr{Priority \\ (High/Med/Low)} &
\Hdr{Space / Safety / Budget Notes} \\

7.17 Aquatic Training and Water Safety Equipment (expanded aquatic logistics) &
T1 → T2 (End) &
Swim becomes consistently gate-relevant; expanded posture supports repeatability and safer execution under variable access. Underwater remains conditional and governance-gated. &
Med &
Space: Moderate. Safety: Surface-only capability readiness; underwater not implied. Budget: Expand only to preserve safe repeatability. \\

\end{longtblr}

\endgroup

\vspace{0.65em}

\subsubsection*{Phase 11 — Pre-Selection Conditioning and Recovery}
\phantomsection

\begingroup
\scriptsize
\setlength{\tabcolsep}{2pt}
\renewcommand{\arraystretch}{1.12}

\begin{longtblr}{
  width=\linewidth,
  rowhead=1,
  colspec = {
    X[l,1.95]
    X[l,0.90]
    X[l,2.40]
    X[l,0.85]
    X[l,2.15]
  },
  hlines,
  vlines,
  cells = {hsep=6pt, vsep=6pt, halign=l, valign=t},
  cell{1-Z}{1-5} = {fg=black},
  row{odd} = {bg=white},
  row{even} = {bg=LightGray},
  row{1} = {bg=StageBlue!15, fg=black, font=\bfseries, halign=l, valign=t},
}
\Hdr{Item or Category Change} &
\Hdr{From Tier → To Tier} &
\Hdr{Why Introduced Here \\ (Training/Test Need)} &
\Hdr{Priority \\ (High/Med/Low)} &
\Hdr{Space / Safety / Budget Notes} \\

7.10 Logistics, Maintenance, Storage, and Sustainment Equipment (maintain Tier 2; prevent last-phase friction) &
T2 → T2 (Operational delta: verification reliability) &
Late phases punish execution friction; sustainment reduces last-minute failures that invalidate test weeks. &
Med &
Space: Moderate. Safety: Prevents illness/skin drift and missed sessions. Budget: Maintain; avoid late upgrades that create novelty. \\

\end{longtblr}

\endgroup

\newpage
\vspace{0.65em}

\subsubsection*{Phase 12 — Final Pre-Selection Conditioning}
\phantomsection

\begingroup
\scriptsize
\setlength{\tabcolsep}{2pt}
\renewcommand{\arraystretch}{1.12}

\begin{longtblr}{
  width=\linewidth,
  rowhead=1,
  colspec = {
    X[l,1.95]
    X[l,0.90]
    X[l,2.40]
    X[l,0.85]
    X[l,2.15]
  },
  hlines,
  vlines,
  cells = {hsep=6pt, vsep=6pt, halign=l, valign=t},
  cell{1-Z}{1-5} = {fg=black},
  row{odd} = {bg=white},
  row{even} = {bg=LightGray},
  row{1} = {bg=StageBlue!15, fg=black, font=\bfseries, halign=l, valign=t},
}
\Hdr{Item or Category Change} &
\Hdr{From Tier → To Tier} &
\Hdr{Why Introduced Here \\ (Training/Test Need)} &
\Hdr{Priority \\ (High/Med/Low)} &
\Hdr{Space / Safety / Budget Notes} \\

7.18 Testing, Timing, Measurement, and Course-Setup Tools (verification redundancy) &
T2 → T3 (End) &
Overbuild for verification reliability; redundancy reduces variance and prevents a single tool failure from invalidating final gates. &
Med &
Space: Low. Safety: Prevents invalid evidence due to tool failure. Budget: Overbuild only where it directly protects admissibility. \\

7.11 Footwear Systems and Foot Care (redundant interface readiness) &
T2 → T3 (End) &
Final verification cannot tolerate foot breakdown; redundancy and sustainment reduce variability and last-minute changes. &
Med &
Space: Low-to-moderate. Safety: Prevents hotspot/blister cascade. Budget: Overbuild only to stabilize; no performance chasing. \\

7.07 Load-Bearing, Rucking, and Carry Systems (redundant fit/config stability) &
T2 → T3 (End) &
Final long-load gates require stable configuration and redundancy to prevent failure due to interface issues. &
Med &
Space: Moderate. Safety: Prevents chafing/fit drift; protects validity. Budget: Overbuild only for reliability. \\

7.12 Recovery, Mobility, and Tissue-Care Implements (redundant recovery reliability) &
T2 → T3 (End) &
Verification blocks require stable recovery capability; redundancy reduces late-phase breakdown. &
Low &
Space: Low. Safety: Supports durability and sleep stability. Budget: Overbuild only if recurring failure patterns exist. \\

7.10 Logistics, Maintenance, Storage, and Sustainment Equipment (verification reliability) &
T2 → T3 (End) &
Sustained high-consequence windows require near-zero friction and high reliability; redundancy prevents avoidable misses. &
Low &
Space: Moderate. Safety: Prevents illness/skin drift. Budget: Only elevate if repeated sustainment failures appear. \\

7.19 Protective and Assistive Equipment (redundant protective posture) &
T2 → T3 (End) &
Final readiness benefits from redundancy where protective failures would derail gates. &
Low &
Space: Low. Safety: Risk containment. Budget: Elevate only when the protective item prevents repeated gate disruption. \\

7.16 Operational Self-Management Systems (redundant planning discipline) &
T2 → T3 (End) &
Verification blocks require strict reslot discipline and no-novelty posture; redundancy reduces drift. &
Low &
Space: Low. Safety: Prevents stacking drift. Budget: Keep minimal unless needed for compliance stability. \\

7.17 Aquatic Training and Water Safety Equipment (verification reliability) &
T2 → T3 (End) &
If \textbf{SWIM} is gate-relevant late-phase, reliability redundancy reduces schedule collapse due to access/tool drift. Underwater remains governance-gated and not implied. &
Low &
Space: Moderate. Safety: Surface-only reliability posture. Budget: Overbuild only if it prevents repeated invalidity. \\

\end{longtblr}

\endgroup

\vspace{0.65em}

\subsubsection*{Phase 13 — Full-Spectrum Readiness Consolidation}
\phantomsection

\begingroup
\scriptsize
\setlength{\tabcolsep}{2pt}
\renewcommand{\arraystretch}{1.12}

\begin{longtblr}{
  width=\linewidth,
  rowhead=1,
  colspec = {
    X[l,1.95]
    X[l,0.90]
    X[l,2.40]
    X[l,0.85]
    X[l,2.15]
  },
  hlines,
  vlines,
  cells = {hsep=6pt, vsep=6pt, halign=l, valign=t},
  cell{1-Z}{1-5} = {fg=black},
  row{odd} = {bg=white},
  row{even} = {bg=LightGray},
  row{1} = {bg=StageBlue!15, fg=black, font=\bfseries, halign=l, valign=t},
}
\Hdr{Item or Category Change} &
\Hdr{From Tier → To Tier} &
\Hdr{Why Introduced Here \\ (Training/Test Need)} &
\Hdr{Priority \\ (High/Med/Low)} &
\Hdr{Space / Safety / Budget Notes} \\

7.18 Testing, Timing, Measurement, and Course-Setup Tools (maintain verification redundancy) &
T3 → T3 (Operational delta: maintain) &
Consolidation assumes verification reliability posture is already achieved; avoid last-minute tool changes that create novelty. &
Med &
Space: Low. Safety: Prevents invalid evidence. Budget: Maintain; no new introductions. \\

7.07 Load-Bearing, Rucking, and Carry Systems (maintain redundancy) &
T3 → T3 (Operational delta: maintain) &
Consolidation maintains stable ruck interface; no novelty near final gates. &
Med &
Space: Moderate. Safety: Prevents interface drift. Budget: Maintain; no new introductions. \\

7.11 Footwear Systems and Foot Care (maintain redundancy) &
T3 → T3 (Operational delta: maintain) &
Final repeatability requires stable foot integrity; do not alter interface near verification. &
Med &
Space: Low-to-moderate. Safety: Prevents blister cascade. Budget: Maintain; no new introductions. \\

7.10 Logistics, Maintenance, Storage, and Sustainment Equipment (maintain redundancy) &
T3 → T3 (Operational delta: maintain) &
Sustainment reliability preserves cadence and prevents avoidable illness/skin drift. &
Low &
Space: Moderate. Safety: Continuity and hygiene. Budget: Maintain; no new introductions. \\

7.19 Protective and Assistive Equipment (maintain redundancy) &
T3 → T3 (Operational delta: maintain) &
Maintain protective posture to avoid preventable breakdown during final proof windows. &
Low &
Space: Low. Safety: Risk containment. Budget: Maintain; no new introductions. \\

\end{longtblr}

\endgroup

\end{landscape}

\end{document}